\chapter{IDEAR}
El presente capítulo desarrolla la tercera etapa del Design Thinking: Idear. Su propósito es generar múltiples soluciones posibles a partir del planteamiento del problema (POV) formulado en el Capítulo 6, sin restricciones prematuras, para después seleccionar aquellas que mejor respondan a las necesidades identificadas durante las fases de Empatizar y Definir (Vargas Márquez et al., 2021).

\section{GENERACIÓN DE IDEAS}
A partir de las preguntas ¿Cómo podríamos...? (Figura 6.6), la generación de ideas se desarrolló mediante dos técnicas complementarias: ¿Y si...? y Lluvia de ideas.

\subsection{¿Y si...?}
La técnica ¿Y si...? tiene como objetivo generar ideas mediante el fomento del pensamiento lateral, explorando escenarios hipotéticos y desafiando suposiciones establecidas (Design Thinking en Español, s.f.-j). Se aplicó en dos rondas, formulando cuatro preguntas por cada uno de los cinco conocimientos (hechos) derivados de la fase de Empatizar, lo que dio lugar a 40 ideas. Entre ellas se encuentran, por ejemplo:

\begin{itemize}
    \item \textbf{Intereses en común:} ¿Y si el emparejamiento fuera como un juego donde el usuario elige ``¿De qué quiero hablar hoy?'' y se conecta con alguien que eligió el mismo tema? ¿Y si los intereses se revelaran progresivamente después de una cantidad de mensajes?
    \item \textbf{Avatar:} ¿Y si el avatar cambiara de expresión según el estado de ánimo que el usuario elige al entrar? ¿Y si el avatar solo se revelara completo después de cierto número de mensajes?
    \item \textbf{Temor a la hostilidad:} ¿Y si, antes de habilitar el chat, ambos usuarios debieran aceptar un breve acuerdo de respeto? ¿Y si existiera un botón de pausa, distinto a bloquear o reportar?
    \item \textbf{Pocos vínculos, pero significativos:} ¿Y si el sistema solo permitiera una conexión activa a la vez? ¿Y si el vínculo se representara como una planta? ¿Y si al finalizar los usuarios pudieran votar de forma anónima si desean volver a hablar, mostrando primero una pantalla de resultados igual para todos para evitar deducir el voto de la otra persona?
    \item \textbf{Mensajes positivos sin métricas:} ¿Y si existiera un banco de frases positivas para cuando el usuario no sabe qué decir? ¿Y si hubiera pequeños gestos no textuales, como un ícono de abrazo o apoyo?
\end{itemize}

\subsection{Lluvia de ideas}
La Lluvia de ideas o \textit{brainstorming} es una técnica que sirve para generar el mayor número posible de ideas, buscando primero la cantidad para aplicar luego técnicas de selección que permitan quedarse con las mejores (Design Thinking en Español, s.f.-d). En el presente proyecto, las ideas más prometedoras de la técnica anterior se desarrollaron como bocetos de concepto (\textit{wireframes}) para cada conocimiento: el emparejamiento por palabra clave de interés y la rotación de temas con aviso; el avatar con expresiones según el estado de ánimo; el acuerdo de respeto previo al chat y el reporte o bloqueo accesible sin interrumpir la experiencia; la opción de reencuentro tras la conversación y el vínculo representado como una planta; y el banco de frases positivas y los mensajes positivos predefinidos para completar.

Varias de las ideas seleccionadas para el prototipo (Tabla \ref{tab:ideas_seleccionadas}) provienen directamente de estas técnicas: la elección de ``¿De qué quiero hablar hoy?'' se transformó en la selección de hasta tres temas antes de cada búsqueda; el avatar que cambia de expresión según el estado de ánimo se incorporó tal como fue planteado; la conexión activa única y la planta dieron forma al ciclo semanal; la votación anónima y la protección del voto de la otra persona se materializaron en la decisión privada al finalizar los siete días; el acuerdo de respeto dio origen a la pantalla ``Un lugar seguro''; y el banco de frases positivas se integró en las ideas para romper el hielo. Algunas ideas se adaptaron al seleccionarse; por ejemplo, la planta que se marchita si no se cuida se transformó en una planta que nunca se marchita.

\section{CRITERIOS DE SELECCIÓN}
Siguiendo a Vargas Márquez et al. (2021), quienes señalan que la creación de múltiples ideas aumenta el potencial de innovación y que es necesario separar la generación de ideas de su evaluación (p. 13), no todas las ideas generadas se llevaron al prototipo. La selección se realizó aplicando tres criterios:

\begin{enumerate}
    \item \textbf{Respuesta a un hallazgo consistente} de las fases de Empatizar y Definir, ya sea un patrón del Mapa de Motivaciones (necesidad de tema en común, 9 de 15 pares; uso de avatar, 9 de 15 pares; entorno seguro o libre de hostilidad, 8 de 15 pares) o un factor identificado en el Capítulo 4.
    \item \textbf{Respaldo en la literatura} revisada en el marco teórico y en el Capítulo 4.
    \item \textbf{Ausencia de carga cognitiva innecesaria}, conforme a la Ley de Hick y la Ley de Tesler (Yablonski, 2020). Las ideas que sumaban opciones a la interfaz sin resolver un dolor central identificado en el Mapa de Empatía se pospusieron para una fase posterior al prototipo actual.
\end{enumerate}

\section{IDEAS SELECCIONADAS}
La Tabla \ref{tab:ideas_seleccionadas} presenta las ideas seleccionadas para el prototipo, organizadas según el conocimiento del cual se derivaron, junto con la pregunta orientadora que las originó y su sustento teórico o empírico.

\begin{longtable}{
  >{\raggedright\arraybackslash}p{0.27\textwidth}
  >{\raggedright\arraybackslash}p{0.35\textwidth}
  >{\raggedright\arraybackslash}p{0.29\textwidth}
}
\caption{Ideas seleccionadas para el prototipo}\label{tab:ideas_seleccionadas}\\
\toprule
\textbf{Pregunta ¿Cómo podríamos...?} & \textbf{Idea seleccionada} & \textbf{Sustento} \\
\midrule
\endfirsthead

\multicolumn{3}{l}{\textit{Tabla \thetable{} (continuación)}} \\[4pt]
\toprule
\textbf{Pregunta ¿Cómo podríamos...?} & \textbf{Idea seleccionada} & \textbf{Sustento} \\
\midrule
\endhead

\midrule
\multicolumn{3}{r}{\textit{Continúa en la siguiente página}} \\
\endfoot

\bottomrule
\endlastfoot

\multicolumn{3}{l}{\textbf{Conocimiento: Compartir gustos o intereses para iniciar y sostener la conversación}} \\[4pt]
¿Cómo podríamos identificar los intereses antes de emparejar?
&
Elección de intereses en el \textit{onboarding}; antes de cada búsqueda se eligen hasta 3 temas, solo entre los propios intereses del usuario.
&
Mapa de Motivaciones (9 de 15 pares); Montoya et al. (2008); Usuarios 1 y 5 \\[6pt]

¿Cómo podríamos emparejar por gustos específicos?
&
Jerarquía de intereses: categoría $\rightarrow$ interés $\rightarrow$ subtema (género, estilo, liga) $\rightarrow$ ejemplo concreto (una película, una banda). El emparejamiento otorga mayor peso a lo más específico.
&
Usuario 4, que solicitó gustos ``específicos y particulares''; Usuario 1, cuyo único recuerdo de conexión positiva surgió de compartir un mismo videojuego \\[6pt]

¿Cómo podríamos mostrar qué tienen en común antes de chatear?
&
Pantalla de presentación que muestra el tema, los intereses y los gustos concretos compartidos entre ambos usuarios.
&
Causa raíz de los 5 ¿Por qué? del Usuario 1: ``no sé de qué hablar'' \\[6pt]

¿Cómo podríamos sugerir temas basados en gustos compartidos?
&
Ideas para romper el hielo: tres sugerencias diarias basadas en lo que ambos comparten; al seleccionarlas se escriben en el campo de texto para su edición, no se envían automáticamente.
&
Usuario 1; la falta de personalización es una barrera común para la participación en intervenciones digitales de bienestar (Borghouts et al., 2021); Ley de Hick (máximo tres opciones) \\[6pt]

\addlinespace[4pt]
\multicolumn{3}{l}{\textbf{Conocimiento: Preferencia por un avatar en vez de una foto real}} \\[4pt]
¿Cómo podríamos permitir que el avatar refleje la identidad real?
&
Avatar personalizable con seis tonos de piel, cinco peinados, seis colores de cabello, y accesorios opcionales (lentes, pecas, lunar, arete), además de ropa y fondo.
&
Usuario 4 (``lentes, pelo largo, lunares, piercings''); Mapa de Motivaciones (9 de 15 pares) \\[6pt]

¿Cómo podríamos hacer el avatar entretenido?
&
Botón ``Sorpréndeme'' que genera una combinación aleatoria de rasgos, con una breve animación de respuesta.
&
Faceta de lo deseable (Morville, 2005, como se citó en Méndez et al., 2014); Efecto de Usabilidad Estética (Yablonski, 2020) \\[6pt]

¿Cómo podríamos compensar la falta de expresiones faciales?
&
El avatar refleja el estado de ánimo del día (ojos, cejas y boca), visible para la otra persona salvo que se desactive.
&
Causa raíz del Usuario 1: ``no puedo ver expresiones, es muy plano''; Picard (1997); Tufail et al. (2025) \\[6pt]

\addlinespace[4pt]
\multicolumn{3}{l}{\textbf{Conocimiento: Temor a la hostilidad o a desconocidos con malas intenciones}} \\[4pt]
¿Cómo podríamos dar reporte y bloqueo sin interrumpir la experiencia?
&
Ícono de escudo dentro de la conversación que lleva a ``Tu bienestar es primero'': opción de salir o reportar con motivo. El motivo no se revela a la otra persona y ambos usuarios no vuelven a ser emparejados.
&
Mapa de Motivaciones (8 de 15 pares); Usuarios 4 y 5; prevención de errores e instrucciones constructivas para recuperarse (Shneiderman et al., 2016) \\[6pt]

¿Cómo podríamos transmitir confianza antes de la primera conversación?
&
Pantalla ``Un lugar seguro'' con acuerdos breves de convivencia y el recordatorio de que se puede salir en cualquier momento.
&
Síntesis de la fase Empatizar (4.5), necesidad (d) \\[6pt]

\addlinespace[4pt]
\multicolumn{3}{l}{\textbf{Conocimiento: Pocos vínculos, pero significativos}} \\[4pt]
¿Cómo podríamos priorizar la calidad sobre la cantidad?
&
Una sola conversación activa a la vez, con duración de siete días y hasta cinco mensajes por persona al día (máximo 500 caracteres cada uno).
&
Burke y Kraut (2016): la comunicación redactada y dirigida personalmente, y no la retroalimentación de un clic, se asocia con mejoras en el bienestar; Saffer (2013). El límite de 500 caracteres atiende la reserva del Usuario 4 \\[6pt]

¿Cómo podríamos hacer que cada mensaje se sienta significativo?
&
Planta que crece solo los días en que ambos usuarios se escriben, con una micro-interacción al avanzar de etapa. Nunca se marchita.
&
Godard y Holtzman (2024): el uso activo se asocia con mayor bienestar ($r = .15$) y apoyo social en línea ($r = .34$) \\[6pt]

¿Cómo podríamos sostener un vínculo esporádico?
&
Decisión privada al finalizar los siete días; si ambos deciden continuar, la planta florece y la conversación permanece abierta.
&
Planteamiento del problema (POV); Regla del Pico-Final (Yablonski, 2020) \\[6pt]

¿Cómo podríamos medir el éxito en bienestar y no en uso?
&
Registro diario del estado de ánimo y una vista privada semanal (``Tu semana''), sin rachas ni notificaciones de enganche.
&
El monitoreo del estado de ánimo cuenta con respaldo empírico y se asocia con mayor autoconciencia emocional (Eisenstadt et al., 2021); elementos que prolongan el uso (Montag et al., 2019) \\[6pt]

\addlinespace[4pt]
\multicolumn{3}{l}{\textbf{Conocimiento: Mensajes positivos, sin 	extit{likes} ni métricas públicas}} \\[4pt]
¿Cómo podríamos sugerir mensajes positivos cuando no se sabe qué decir?
&
Frases de validación integradas en las ideas sugeridas (ej. ``Gracias por quedarte en la conversación'').
&
Pedido explícito del Usuario 1; la gratitud fortalece los vínculos (Algoe, 2012); Clark et al. (2018) \\[6pt]

¿Cómo podríamos evitar métricas de comparación?
&
Ausencia total de contadores públicos: el Jardín muestra plantas, no cifras ni rankings.
&
Factor 4 del Capítulo 4 (Montag et al., 2019; Verduyn et al., 2017); límites del alcance del proyecto \\[6pt]

¿Cómo podríamos hacer que la app acompañe en un mal día?
&
Respuestas empáticas según el ánimo registrado y un ejercicio de respiración guiada de 25 segundos cuando el ánimo es bajo.
&
Eisenstadt et al. (2021): respaldo a los ejercicios guiados y al monitoreo del estado de ánimo en aplicaciones de bienestar \\

\end{longtable}

\section{IDEAS DESCARTADAS O POSPUESTAS}
No todas las ideas generadas durante la fase de Idear se incorporaron al prototipo. La Tabla \ref{tab:ideas_descartadas} sintetiza las ideas descartadas o pospuestas, junto con la razón de su exclusión, en coherencia con la recomendación de separar la generación de ideas de su evaluación (Vargas Márquez et al., 2021, p. 13).

\begin{longtable}{
  >{\raggedright\arraybackslash}p{0.28\textwidth}
  >{\raggedright\arraybackslash}p{0.17\textwidth}
  >{\raggedright\arraybackslash}p{0.47\textwidth}
}
\caption{Ideas descartadas o pospuestas}\label{tab:ideas_descartadas}\\
\toprule
\textbf{Idea} & \textbf{Decisión} & \textbf{Razón} \\
\midrule
\endfirsthead

\multicolumn{3}{l}{\textit{Tabla \thetable{} (continuación)}} \\[4pt]
\toprule
\textbf{Idea} & \textbf{Decisión} & \textbf{Razón} \\
\midrule
\endhead

\midrule
\multicolumn{3}{r}{\textit{Continúa en la siguiente página}} \\
\endfoot

\bottomrule
\endlastfoot

``Mundo de conexiones'' con estadísticas globales de usuarios & Descartada & Introducía métricas públicas y comparación social, contrario al factor 4 identificado en el Capítulo 4 \\[6pt]
Notificaciones \textit{push} y rachas diarias (\textit{streaks}) & Descartada & Son elementos que prolongan el tiempo de uso (Montag et al., 2019) y podrían fomentar un patrón de uso problemático, asociado a mayor soledad (Feng et al., 2026) \\[6pt]
Reacciones variadas a los mensajes (solicitado por el Usuario 1) & Pospuesta & Suma opciones de interacción en cada mensaje (Ley de Hick, Yablonski, 2020); candidata a incorporarse tras la prueba piloto \\[6pt]
Mensajes largos tipo \textit{story time} (solicitado por el Usuario 4) & Atendida parcialmente & Se amplió el límite a 500 caracteres por mensaje, sin introducir un formato especial que rompa la lógica de micro-interacción \\[6pt]
Debates de opinión grupales (solicitado por el Usuario 4) & Descartada & Excede el alcance definido en 1.5, que excluye la comunicación grupal \\[6pt]
Avatar en modelado 3D & Reemplazada & Sustituida por un avatar 2D vectorial (SVG), más liviano y adecuado para personalización en tiempo real en dispositivos móviles \\[6pt]
Lista fija de seis temas para la búsqueda de pareja & Reemplazada & Mostraba temas no elegidos por el usuario; se sustituyó por una búsqueda basada en sus propios intereses \\

\end{longtable}

\section{SÍNTESIS DE LA FASE IDEAR}
Las ideas seleccionadas y presentadas en la Tabla \ref{tab:ideas_seleccionadas} constituyen el conjunto de decisiones de diseño que se materializan en el prototipo funcional descrito en el Capítulo 8. Cada una de ellas responde directamente a un hallazgo de las fases de Empatizar y Definir, y cuenta con respaldo en la literatura revisada en los Capítulos 2 y 4, en coherencia con el enfoque iterativo y centrado en el usuario del Design Thinking (Vargas Márquez et al., 2021).